\documentclass[journal]{IEEEtran}
\usepackage[utf8]{inputenc}
\usepackage[T1]{fontenc}
\usepackage{amsmath,amssymb,amsfonts}
\usepackage{algorithm}
\usepackage{algorithmic}
\usepackage{graphicx}
\usepackage{textcomp}
\usepackage{xcolor}
\usepackage{hyperref}
\usepackage{booktabs}
\usepackage{listings}
\usepackage{bm}

\begin{document}

\title{LabelScope: A Claim-Capping Framework for Auditing Operational Labels in Physiology-Based Human-Machine Systems}

\author{Kotaro Furukawa}

\maketitle

\begin{abstract}
Physiology-based human-machine systems (HMS) often train predictive models on operational labels such as questionnaire scores, task conditions, behavioral errors, or affective annotations. However, high predictive accuracy does not by itself establish that the target label can support the HMS claim made from the model. A label may be temporally sparse, behaviorally confounded, or structurally unsupported in physiological feature space. In such cases, model performance may reflect the learnability of an operational proxy rather than evidence for a physiology-grounded user-state claim. This paper introduces LabelScope, a conservative pre-modeling audit framework for claim-relative state-label validity in physiology-based HMS. LabelScope audits whether an operational label can support a specified HMS claim under a given modeling resolution, proxy interpretation, feature representation, and deployment context. The framework consists of Resolution, Proxy, Structure, and Claim audits, and assigns a capped claim level $L_{Final} \in \{0,1,2,3\}$ using a conservative claim-capping rule. We illustrate LabelScope through worked audit demonstrations using two public physiological datasets: SWELL-KW and WESAD. These demonstrations illustrate how LabelScope caps unsupported claims arising from resolution mismatch and behavioral proxy confounding. LabelScope does not determine whether a psychological state exists or whether a construct is invalid; rather, it specifies the strongest defensible HMS claim supported by the tested operational label and pipeline. The framework provides a reporting discipline for interpreting predictive accuracy together with audited claim levels in physiology-based HMS research.
\end{abstract}

\begin{IEEEkeywords}
Physiology-based HMS, State-Label Validity, Claim-Capping, Auditing, Human-Machine Systems, Affective Computing, Behavioral Proxy.
\end{IEEEkeywords}
\section{Introduction}

\IEEEPARstart{T}{he} integration of physiological sensing into human-machine systems (HMS) represents one of the most promising frontiers in the design of intelligent interaction. By monitoring signals such as heart rate variability (HRV), electrodermal activity (EDA), and electroencephalography (EEG), researchers aim to build systems that can sense and respond to a user's internal state in real-time. These systems, often categorized under affective computing \cite{picard1997} or cognitive load monitoring \cite{martinez2017}, promise to revolutionize a wide range of domains, from healthcare and industrial safety to personalized learning and entertainment. The fundamental promise of physiological HMS is the ability to bypass the "bottleneck" of explicit user input, allowing systems to adapt autonomously to a user's stress, fatigue, or engagement \cite{parasuraman2000}.

Over the past decade, the rapid advancement of machine learning algorithms and the increasing availability of public physiological datasets have significantly lowered the barrier to entry for building these systems. Researchers can now extract thousands of features from physiological signals and train complex deep learning models to predict a variety of "ground truth" labels. However, as the field has matured, a critical methodological gap has emerged: the tendency to equate predictive accuracy with the validity of the underlying user-state claim. In many HMS studies, a high-accuracy model trained on an operational label---such as a questionnaire score or a specific task condition---is frequently treated as definitive evidence that the system has successfully "recognized" or "detected" the intended psychological state \cite{calvo2010}. 

We argue that this focus on predictive accuracy creates a \textbf{methodological risk} in the interpretation of results. High accuracy measures only the learnability of an operational target within a specific dataset and feature representation. It does not, by itself, justify the scope of the HMS claim made from that model. The central problem addressed in this paper is that operational labels in physiology-based HMS are often underdetermined by the physiological evidence. A model may achieve high accuracy by exploiting behavioral confounding, pseudo-replication artifacts, or dataset-specific noise, rather than by capturing a robust physiological signature of the intended psychological state. 

This underdetermination creates a profound risk for Human-Machine Systems (HMS). In a closed-loop HMS, the system's action policy is directly coupled to the inferred user state. If a system triggers an intervention based on a confounded label---for example, an adaptive learning environment incorrectly lowering task difficulty because it mistook a student's physical posture shift (a Level 2 behavioral proxy) for cognitive overload (a Level 3 internal state)---the resulting action disrupts the user's workflow. Such misdirected interventions not only degrade task performance but also rapidly erode the user's trust in the system's autonomy \cite{parasuraman2000}. To prevent these silent failures in human-machine interaction, state-label validity must be treated as a systematic pre-modeling audit discipline before predictive performance is allowed to govern system behavior.

In this paper, we introduce \textbf{LabelScope}. This paper proposes a claim-relative audit framework that caps the strongest defensible HMS inference supported by an operational physiological label before predictive accuracy is interpreted. By treating validity as a claim-relative property, LabelScope allows researchers to formally specify the strongest defensible interpretation of their models, whether as a coarse behavioral proxy or a fine-grained internal state monitor.
\section{Related Work}

\subsection{Physiological HMS and the Affective Computing Heritage}
Physiology-based HMS are deeply rooted in the foundational goals of affective computing. Since Picard's seminal work in 1997 \cite{picard1997}, the field has sought to build computational systems that can relate to, arise from, or influence human emotions and internal states. This heritage has provided the technical scaffolding for modern HMS, including signal processing techniques for heart rate variability (HRV), skin conductance \cite{boucsein2012}, and brain-computer interfaces. Over the past two decades, the focus has shifted from basic emotion recognition to the sensing of complex states such as cognitive load and stress. 

Surveys of the field, such as those by Calvo and D'Mello \cite{calvo2010} and Bota et al. \cite{bota2019}, highlight the increasing sophistication of physiological sensing pipelines. However, as sensor technology and modeling techniques have advanced, a new phase of methodological refinement has become necessary. While initial studies established the foundational links between physiology and psychology, the deployment of these models in adaptive HMS requires a more granular understanding of state-label supportability. The challenge of establishing a ``ground truth'' for unobservable states remains a central methodological question, one that now requires explicit auditing of the interpretive scope of our models.

\subsection{Adaptive Automation and Human-Machine Teaming}
The ultimate goal of physiological HMS is to enable adaptive automation, where the system dynamically adjusts its intervention (e.g., task offloading, alert suppression) based on the operator's cognitive or affective state \cite{parasuraman2000}. A critical challenge in human-machine teaming is trust calibration \cite{lee2004}. When an adaptive system intervenes based on a spurious state inference, it degrades the operator's trust and situational awareness \cite{wickens2008}. LabelScope addresses this by providing a pre-intervention audit, ensuring that the system's action policy is grounded in a defensible physiological claim rather than an unverified operational proxy.

\subsection{Construct Validity and the Ground Truth Problem}
The problem of ground truth in affective and physiological computing is notoriously difficult \cite{calvo2010}. Systems typically rely on subjective self-reports or task-based objective conditions. However, this relates to the broader theory of Construct Validity \cite{messick1989}, which emphasizes that validity is a property of the \textit{interpretations} of scores, not the scores themselves. While traditional frameworks evaluate validity conceptually, LabelScope operationalizes construct validity as a computable, claim-relative metric. We treat the operational label as a weak or proxy label \cite{zhou2018} that must be mathematically audited for its supportability at a specific resolution.

\subsection{Leakage and Confounding in Physiological ML}
Recent literature has highlighted severe methodological flaws in physiological ML, particularly data leakage through pseudo-replication (where overlapping windows from the same label observation inflate accuracy) \cite{pseudo2015} and behavioral confounding (where physiological sensors act as high-variance motion detectors) \cite{chiossi2024}. LabelScope is designed to integrate these known failure modes into an explicit, deterministic auditing pipeline that outputs a \textbf{claim cap}. While general ML accountability tools like Model Cards \cite{mitchell2019} and Datasheets \cite{gebru2021} provide descriptive metadata, LabelScope provides a prescriptive algorithmic bound on the physiological interpretability of the model.
\section{The LabelScope Framework}

\subsection{Mathematical Formalization}
Let $X \in \mathbb{R}^{N \times D}$ be a physiological feature matrix. Let $y \in \mathbb{R}^N$ be the operational label, and $C_B \in \mathbb{R}^{N \times K}$ be a matrix of $K$ behavioral covariates. LabelScope evaluates the quadruple $(X, y, C_B, S)$ and returns a final claim level $L_{Final} \in \{0, 1, 2, 3\}$.

\subsection{Audit Module 1: Resolution Audit (RA)}
Resolution Audit (RA) addresses the risk of \textbf{pseudo-replication}, which is a major source of inflated accuracy in HMS research \cite{pseudo2015}. Pseudo-replication occurs when a model is trained on multiple samples that are not independent but share the same label observation. 

We define the replication ratio $R_P$ as:
\begin{equation}
R_P = \frac{N_{samples}}{N_{independent\_observations}}
\end{equation}
If $R_P$ exceeds the RA threshold $\tau_{RA}$, the claim is capped at Level 0. This ensures that the modeling resolution matches the information frequency of the labels.

\subsection{Audit Module 2: Proxy Audit (PA)}
Proxy Audit (PA) detects behavioral confounding. Physiological signals are highly sensitive to physical activity. If an operational label is strongly correlated with these behaviors, the system may be acting as a behavioral proxy detector. We define the proxy index $\rho_{LC}$ as:
\begin{equation}
\rho_{LC} = \max_{k \in \{1 \dots K\}} |\mathrm{Corr}(y, c_k)|
\end{equation}
If $\rho_{LC}$ exceeds the PA threshold $\tau_{PA}$, the claim is capped at Level 2. This requires including non-physiological sensors (e.g., accelerometers) in the auditing pipeline. It is critical to note that passing the PA only clears the system of \textit{observed} covariates ($C_B$). The PA cannot guarantee the absence of unobserved behavioral or environmental confounders (e.g., unreported postural shifts, room temperature changes). Therefore, a passed PA provides a defensive lower bound on risk, not an absolute guarantee of state purity.

\subsection{Audit Module 3: Structure Audit (SA)}
Structure Audit (SA) evaluates whether the operational label corresponds to a non-random local structure in the physiological feature space. We define the test statistic $T_{obs}$ as the $k$-Nearest Neighbors (k-NN) label coherence:
\begin{equation}
T_{obs} = \frac{1}{N} \sum_{i=1}^{N} \mathbf{1}\left(y_i = \text{mode}(\{y_j\}_{j \in \mathcal{N}_k(x_i)})\right)
\end{equation}
where $\mathcal{N}_k(x_i)$ is the set of $k$ nearest neighbors of $x_i$ using Euclidean distance after subject-wise z-score scaling. To account for temporal autocorrelation and avoid artificially inflated coherence, the permutation test uses \textbf{block-shuffling} rather than naive sample shuffling. The label time-series is divided into contiguous task blocks, and the blocks are permuted within each subject before calculating the surrogate coherence $T_{shuf}^{(b)}$. The empirical p-value is calculated as:
\begin{equation}
\hat{p}_{shuf} = \frac{1 + \sum_{b=1}^{B} \mathbf{1}\left(T_{shuf}^{(b)} \ge T_{obs}\right)}{B + 1}
\end{equation}
If $\hat{p}_{shuf}$ fails to achieve statistical significance at level $\alpha$, the claim is capped at Level 1. This audit protects against ``accuracy from noise'' in high-dimensional space.

\subsection{Audit Module 4: Claim Audit and Relativity}
The final claim level $L_{Final}$ is determined by the minimum cap imposed by RA, PA, and SA. Table \ref{tab:claim_levels} explicitly defines the cap conditions, permitted HMS claims, and prohibited interpretations for each level. 

\begin{table*}[t]
\caption{LabelScope Claim Level Definitions and Interpretative Boundaries}
\label{tab:claim_levels}
\centering
\begin{tabular}{@{}p{1.2cm}p{3cm}p{3.5cm}p{4cm}p{4cm}@{}}
\toprule
\textbf{Level} & \textbf{Designation} & \textbf{Cap Condition} & \textbf{Permitted HMS Claim} & \textbf{Prohibited Interpretation} \\
\midrule
\textbf{Level 0} & Unsupported Resolution & $R_P > \tau_{RA}$ (RA Failure) & None at the intended resolution. Must aggregate data. & Real-time or continuous state detection. \\
\textbf{Level 1} & Dataset/Noise Artifact & $\hat{p}_{shuf} \ge \alpha$ (SA Failure) & Recognition of specific dataset context or random fluctuations. & Any generalization to psychological constructs. \\
\textbf{Level 2} & Behavioral Proxy & $\rho_{LC} > \tau_{PA}$ (PA Failure) & Detection of physical movement, task actions, or behavior. & Direct measurement of internal cognitive/affective state. \\
\textbf{Level 3} & Candidate Valid State & Passes all audits & Strong evidence for physiology-grounded state detection. & Absolute ground truth (unobserved confounders may exist). \\
\bottomrule
\end{tabular}
\end{table*}

It is critical to emphasize the \textbf{relativity of the claim-cap}: the audit outcome is strictly relative to the specifically defined operational label and feature resolution. If a researcher alters the label definition (e.g., aggregating from a ``minute-level'' state to a ``block-level'' state), the pseudo-replication ratio ($R_P$) and the behavioral correlation ($\rho_{LC}$) fundamentally change. Therefore, the audit pipeline must be re-run, potentially resulting in a different cap. The cap is not an absolute property of the underlying psychological construct, but a bounded property of the specific label-feature pair evaluated.

The thresholds $\tau_{RA}$, $\tau_{PA}$, and $\alpha$ are configurable defaults rather than universal constants; their appropriate values are domain-dependent. In the demonstrations that follow, we set $\tau_{RA} = 5.0$, $\tau_{PA} = 0.30$, and $\alpha = 0.05$ as conservative starting points, and discuss their sensitivity in Section~\ref{sec:discussion}.

\section{Signal Processing and Feature Engineering}

\subsection{Pre-Modeling Signal Processing}
Accurate auditing requires a robust signal processing foundation. Signals were first resampled to a common processing rate (e.g., 256 Hz for ECG, 4 Hz for EDA) where necessary to ensure consistent feature extraction. The original dataset-specific hardware acquisition rates are reported in Table \ref{tab:sensors}. In this work, we utilize a multi-modal processing pipeline implemented using \textbf{NeuroKit2} \cite{neurokit2021}. ECG signals are processed using standard bandpass filters. R-peaks are detected and validated to derive HRV features.

EDA signals are processed to separate the tonic and phasic components. We utilize the cvxEDA algorithm \cite{neurokit2021}, which frames the decomposition as a convex optimization problem. These features are critical for ensuring that the Structure Audit reflects state-based variations rather than sensor noise.

\subsection{Feature Extraction and Normalization}
We extracted a comprehensive set of features, including time-domain and frequency-domain HRV. All features were normalized using subject-wise z-score normalization to account for individual differences \cite{cacioppo2007}. This is essential for the Structure Audit to work across subjects.
\section{Experimental Demonstrations}

To illustrate the utility of the LabelScope framework, we applied it to two public physiological datasets: SWELL-KW and WESAD. For each dataset, we defined an intended HMS claim representing a typical research goal and applied LabelScope to audit the supportability of the operational labels under the demonstration thresholds $\tau_{RA} = 5.0$, $\tau_{PA} = 0.30$, $\alpha = 0.05$, $k = 50$, $B = 1000$, and block length $L = 60$\,s.

\subsection{SWELL-KW: Auditing Mental Effort in Knowledge Work}
The SWELL-KW dataset \cite{swell2014} contains physiological signals from 25 participants performing knowledge work (e.g., researching, writing, and presentations) under varying levels of stress and workload. The primary operational label used in many studies is the subjective "Mental Effort" score, collected at the end of each task block (~15--20 minutes).

\subsubsection{Resolution Audit (RA)}
The intended HMS claim was minute-level monitoring of mental effort for adaptive task assistance. However, the mental effort labels were only collected every 15 minutes. For a 1-minute modeling resolution, the replication ratio $R_P$ was 15. With only 75 independent label observations across the entire dataset, $R_P$ exceeded the conservative threshold for fine-grained inference (Table \ref{tab:swell_ra}).

\begin{table*}[t]
\caption{SWELL-KW Resolution Audit (RA) Parameters and Calculation}
\label{tab:swell_ra}
\centering
\begin{tabular}{@{}llcccccc@{}}
\toprule
\textbf{Dataset} & \textbf{Intended Claim} & \textbf{Label Interval} & \textbf{Window Size} & \textbf{Indep. Labels} & \textbf{Samples} & \textbf{$R_P$} & \textbf{Final Cap} \\
\midrule
SWELL-KW & Minute-level Effort & 15 min & 1 min & 75 & 1125 & 15.0 & \textbf{Level 0} \\
\bottomrule
\end{tabular}
\end{table*}

\textbf{Result}: The claim was \textbf{Capped at Level 0}.
\textbf{Interpretation}: While high accuracy has been reported for minute-level prediction in this dataset, LabelScope identifies that these models are learning block-level idiosyncrasies rather than genuine minute-level state transitions. The claim of "minute-level monitoring" is unsupported by the sparse data.

\subsubsection{Proxy Audit (PA)}
We audited the mental effort label against typing volume and mouse activity. While some correlation existed during specific task phases, it was not the primary bottleneck compared to the resolution mismatch.

\subsection{WESAD: Auditing Stress vs. Physical Activity}
The WESAD (Wearable Stress and Affect Detection) dataset \cite{wesad2018} contains multi-modal data from 15 subjects during laboratory induction of stress and affect. The dataset includes high-resolution physiological signals (ECG, EDA, Respiration) and 3D acceleration.

\subsubsection{Proxy Audit (PA)}
The intended claim was continuous sensing of internal stress. However, many of the stress-induction tasks (e.g., public speaking, mental arithmetic) were associated with changes in posture and physical activity. We audited the binary stress vs. non-stress label against the accelerometer-based Signal Magnitude Area (SMA) derived from the wrist-worn Empatica E4. We explicitly defined the `non-stress` condition by pooling the `baseline` and `amusement` protocol segments. The SMA was calculated continuously over a 60-second window ($T=60$) using the formula: $SMA = \frac{1}{T} \int_0^T (|a_x(t)| + |a_y(t)| + |a_z(t)|) dt$.

\textbf{Result}: The Proxy Audit was calculated subject-wise using continuous 60-second sliding windows with 50\% overlap, yielding approximately 60-90 windows per subject. Because the relationship between physiological activation and movement is non-linear but monotonic, we utilized Spearman's rank correlation. Table \ref{tab:wesad_pa} details the subject-wise $\rho_{LC}$ values and corresponding permutation p-values. Across the 15 subjects, the mean proxy index was $\rho_{LC} = 0.65$ ($\sigma = 0.12$, 95\% CI $[0.59, 0.71]$), indicating a strong material association between the label and physical movement.

\begin{table}[h]
\caption{WESAD Subject-wise Proxy Index ($\rho_{LC}$)}
\label{tab:wesad_pa}
\centering
\begin{tabular}{@{}lcc|lcc@{}}
\toprule
\textbf{Subject} & \textbf{$\rho_{LC}$} & \textbf{p-value} & \textbf{Subject} & \textbf{$\rho_{LC}$} & \textbf{p-value} \\
\midrule
S2  & 0.68 & $<0.01$ & S10 & 0.71 & $<0.01$ \\
S3  & 0.62 & $<0.01$ & S11 & 0.64 & $<0.01$ \\
S4  & 0.75 & $<0.01$ & S13 & 0.58 & $<0.01$ \\
S5  & 0.59 & $<0.01$ & S14 & 0.66 & $<0.01$ \\
S6  & 0.81 & $<0.01$ & S15 & 0.49 & $<0.05$ \\
S7  & 0.61 & $<0.01$ & S16 & 0.73 & $<0.01$ \\
S8  & 0.55 & $<0.01$ & S17 & 0.69 & $<0.01$ \\
S9  & 0.65 & $<0.01$ & \textbf{Mean} & \textbf{0.65} & - \\
\bottomrule
\end{tabular}
\end{table}

\textbf{Interpretation}: The claim was \textbf{Capped at Level 2}. The audited label cannot support a claim of movement-independent internal stress sensing; it should instead be reported as a task/protocol-linked stress proxy under the tested behavioral covariates.

\subsubsection{Structure Audit (SA)}
To demonstrate the full sequential audit, we evaluated the WESAD labels under the Structure Audit. Despite the strong behavioral confounding identified by the PA, the stress condition induces massive sympathetic activation (e.g., SCL increases, HRV decreases) that creates a highly structured neighborhood in the physiological feature space.

\textbf{Result}: The empirical p-value $\hat{p}_{shuf}$ was significant ($p < 0.05$), meaning the ``stress vs.\ non-stress'' labels passed the SA and avoided a Level 1 cap. Table \ref{tab:wesad_sa} summarizes the SA parameters and outcome.

\begin{table}[h]
\caption{WESAD Structure Audit (SA) Parameters and Outcome}
\label{tab:wesad_sa}
\centering
\begin{tabular}{@{}lcccccc@{}}
\toprule
\textbf{N win.} & \textbf{D feat.} & \textbf{k} & \textbf{B} & \textbf{$T_{obs}$} & \textbf{$\hat{p}_{shuf}$} & \textbf{SA cap} \\
\midrule
1110 & 19 & 50 & 1000 & 0.74 & 0.001 & None (pass) \\
\bottomrule
\end{tabular}
\end{table}

\textbf{Interpretation}: While the physiological structure is non-random (passing SA), the claim is still bounded at Level 2 by the PA. This demonstrates the necessity of the multi-module framework: passing SA establishes physiological coherence, but only PA can detect when that coherence is driven by concurrent physical activity.

\subsection{Summary of Audit Results}
The demonstrations illustrate how LabelScope identifies diverse failure modes across different datasets.
\begin{itemize}
    \item \textbf{SWELL-KW}: Failure due to resolution mismatch (RA).
    \item \textbf{WESAD}: Failure due to behavioral confounding (PA).
\end{itemize}
In each case, LabelScope provided a clear "cap" on the defensible HMS claim, ensuring that predictive accuracy is not misinterpreted as evidence for an unsupported state-label relationship.
\section{Audit Implementation Details}

\subsection{Algorithmic Procedures}
The core logic of the LabelScope audits is formalized in Algorithms 1 and 2. These procedures ensure that the audit results are reproducible and independent of the specific machine learning model used in the downstream task. For the experimental demonstrations, the following hyperparameters were utilized: k-NN neighborhood size $k=50$, number of permutations $B=1000$, and block-shuffling length $L=60$ seconds.

\begin{algorithm}
\caption{Resolution Audit (RA) Procedure}
\begin{algorithmic}[1]
\REQUIRE Dataset $D$, Intended Resolution $\Delta t$, Independent Observation Frequency $f_{ind}$
\ENSURE Resolution Audit Level $L_{RA}$
\STATE $N_{samples} \gets \text{count samples at } \Delta t$
\STATE $N_{ind} \gets \text{count samples at } f_{ind}$
\STATE $R_P \gets N_{samples} / N_{ind}$
\IF{$R_P > \tau_{RA}$}
    \STATE $L_{RA} \gets 0$
\ELSE
    \STATE $L_{RA} \gets 3$ \COMMENT{No resolution cap imposed}
\ENDIF
\RETURN $L_{RA}$
\end{algorithmic}
\end{algorithm}

\begin{algorithm}
\caption{Structure Audit (SA) Permutation Test}
\begin{algorithmic}[1]
\REQUIRE Feature Matrix $X$, Label Vector $y$, Block Length $L$, Permutations $B$
\ENSURE Empirical p-value $\hat{p}_{shuf}$
\STATE $T_{obs} \gets \text{calculate\_coherence}(X, y)$
\STATE $count \gets 0$
\FOR{$b = 1$ \TO $B$}
    \STATE $y_{shuf} \gets \text{block\_shuffle}(y, L)$
    \STATE $T_{shuf}^{(b)} \gets \text{calculate\_coherence}(X, y_{shuf})$
    \IF{$T_{shuf}^{(b)} \ge T_{obs}$}
        \STATE $count \gets count + 1$
    \ENDIF
\ENDFOR
\STATE $\hat{p}_{shuf} \gets (1 + count) / (B + 1)$
\RETURN $\hat{p}_{shuf}$
\end{algorithmic}
\end{algorithm}

\subsection{Sensor and Feature Specifications}
To ensure reproducibility, we provide the detailed sensor specifications and feature engineering parameters used in the demonstrations (Table \ref{tab:sensors} and \ref{tab:extraction}).

\begin{table}[h]
\caption{Hardware Acquisition Specifications Across Evaluated Datasets}
\label{tab:sensors}
\centering
\begin{tabular}{@{}lll@{}}
\toprule
Dataset & Target Signal & Device \& Sample Rate \\
\midrule
\textbf{WESAD} & Chest ECG, EDA, RESP & RespiBAN Professional (700 Hz) \\
               & Wrist BVP, EDA, ACC  & Empatica E4 (64 Hz / 4 Hz) \\
\midrule
\textbf{SWELL-KW}& Chest ECG, EDA     & TMSi Mobi8 (2048 Hz) \\
\bottomrule
\end{tabular}
\end{table}

\begin{table}[h]
\caption{Feature Extraction Parameters}
\label{tab:extraction}
\centering
\begin{tabular}{@{}lll@{}}
\toprule
Parameter & Value & Description \\
\midrule
Window Size & 60 s & Sliding window duration \\
Overlap & 50 \% & Step size for extraction \\
Bandpass Filter & 0.5--40 Hz & ECG noise removal \\
EDA Threshold & 0.01 $\mu$S & Peak detection floor \\
\bottomrule
\end{tabular}
\end{table}
\section{Discussion: Implications for HMS Design and Policy}
\label{sec:discussion}

\subsection{From Accuracy-Centric to Claim-Centric Reporting}
The primary shift proposed by LabelScope is a transition from accuracy-centric reporting to claim-centric auditing. In current HMS research, a model with 90\% accuracy is almost universally treated as superior to one with 80\% accuracy. However, as our demonstrations show, high accuracy is meaningless if the underlying claim is unsupported. A Level 0 model with 95\% accuracy (driven by pseudo-replication) or a Level 2 model with 90\% accuracy (driven by behavioral confounding) should not be interpreted as evidence for a successful physiology-based HMS.

LabelScope requires researchers to report their predictive performance \textit{within the context of an audited claim level}. This ensures that engineering convenience (e.g., using sparse labels for high-frequency modeling) is not mistaken for support for an HMS claim. We advocate that future HMS research report claim supportability levels alongside traditional performance metrics.

\subsection{LabelScope as an Engineering Safeguard}
Beyond its role as a reporting discipline, LabelScope serves as an engineering safeguard during the HMS design process. By identifying failure modes such as resolution mismatch or behavioral confounding \textit{before} the system is deployed, designers can make informed decisions to improve the system's reliability.
For example, if a "stress monitor" is capped at Level 2 due to movement confounding, the designer can either:
\begin{enumerate}
    \item \textbf{Refine the Hardware}: Add additional sensors to explicitly control for behavioral covariates.
    \item \textbf{Refine the Label}: Collect higher-resolution labels that are less confounded with activity.
    \item \textbf{Refine the Claim}: Explicitly market and deploy the system as a "task-behavior monitor" rather than a pure "physiological state monitor."
\end{enumerate}
This proactive auditing reduces the risk of misdirected interventions and improves the overall safety and trust in adaptive human-machine systems.



\subsection{Limitations and Future Work}
While LabelScope provides a rigorous foundation, it is not a "validity oracle." Its results are sensitive to the choice of features, neighborhood definitions, and audit thresholds. Two critical limitations must be highlighted:

\begin{enumerate}
    \item \textbf{Absence of Closed-Loop Evaluation}: The current study evaluates LabelScope only through retrospective audits of public datasets. It does not evaluate the framework in a closed-loop HMS or with human participants. Determining how audit-triggered claim caps affect actual system behavior, user experience, and task performance requires rigorous user studies in future deployments.
    \item \textbf{Epistemological Limits of the Proxy Audit}: The PA only audits against \textit{observed} covariates. In unrestricted environments, physiological signals can be confounded by unobserved factors (e.g., room temperature, undocumented postural shifts, caffeine intake). A ``Level 3'' claim indicates that the state is unconfounded by measured behaviors, but it cannot guarantee the absence of unmeasured confounders. Future work should integrate causal inference techniques to bound the impact of unobserved variables.
    \item \textbf{Threshold Standardization}: The default values of $\tau_{RA} = 5.0$ and $\tau_{PA} = 0.30$ were chosen conservatively. Establishing domain-specific thresholds through large-scale meta-analysis is necessary to improve precision.
\end{enumerate}

\section{Conclusion}
We have presented LabelScope, a conservative framework for auditing state-label validity in physiology-based HMS. By shifting the focus from predictive accuracy to claim supportability, we provide a reporting discipline that can improve the reliability and reproducibility of physiological computing. As HMS are deployed in increasingly consequential environments, systematic auditing of claim supportability will be an important step toward building human-centered systems with a well-characterized scope of inference.
\appendices
\section{Representative Feature Definitions}

Table \ref{tab:features} lists representative physiological features used in the Structure Audit (SA). All features were extracted using a 60-second sliding window with 50\% overlap. The full feature set ($D = 19$) consists of the HRV time-domain, HRV frequency-domain, HRV nonlinear, EDA tonic, EDA phasic, and accelerometer magnitude domains listed below.

\begin{table*}[!ht]
\caption{Representative Extracted Physiological Features ($D = 19$ total)}
\label{tab:features}
\centering
\begin{tabular}{@{}lll@{}}
\toprule
Signal Domain & Feature Name & Mathematical Definition / Description \\
\midrule
HRV (Time) & Mean RR & Mean of the intervals between successive R-peaks. \\
HRV (Time) & SDNN & Standard deviation of the NN intervals. \\
HRV (Time) & RMSSD & Root mean square of successive RR interval differences. \\
HRV (Time) & pNN50 & Percentage of successive RR intervals that differ by more than 50 ms. \\
HRV (Freq) & VLF Power & Power in the very low frequency range (0.0033--0.04 Hz). \\
HRV (Freq) & LF Power & Power in the low frequency range (0.04--0.15 Hz). \\
HRV (Freq) & HF Power & Power in the high frequency range (0.15--0.4 Hz). \\
HRV (Freq) & LF/HF Ratio & Ratio of LF to HF power, indicating autonomic balance. \\
HRV (Non-lin) & SD1 & Short-term variability from the Poincare plot. \\
HRV (Non-lin) & SD2 & Long-term variability from the Poincare plot. \\
HRV (Non-lin) & Sample Entropy & Measure of the complexity/regularity of the RR time series. \\
EDA (Tonic) & SCL Mean & Mean level of the tonic skin conductance component. \\
EDA (Tonic) & SCL SD & Standard deviation of the tonic skin conductance component. \\
EDA (Phasic) & SCR Peak Count & Number of phasic peaks detected in the window. \\
EDA (Phasic) & SCR Amplitude & Mean amplitude of the detected phasic peaks. \\
EDA (Phasic) & SCR AUC & Area under the curve of the phasic component. \\
ACC (Behav) & SMA & $\frac{1}{T}\int_0^T(|a_x|+|a_y|+|a_z|)dt$; gravity-subtracted per-axis before summing. \\
ACC (Behav) & Mean Accel & Average per-axis acceleration magnitude across $x$, $y$, $z$. \\
\bottomrule
\end{tabular}
\end{table*}

\bibliographystyle{IEEEtran}
\begin{thebibliography}{30}

% Foundations of Affective Computing
\bibitem{picard1997} R. W. Picard, \emph{Affective Computing}. MIT Press, 1997.
\bibitem{calvo2010} R. A. Calvo and S. D'Mello, ``Affect detection: An interdisciplinary review of models, methods, and their applications,'' \emph{IEEE Trans. Affective Comp.}, vol. 1, no. 1, pp. 18--37, 2010.
\bibitem{bota2019} P. Bota et al., ``A review of emotion recognition from physiological signals using machine learning,'' \emph{Sensors}, vol. 19, no. 19, 2019.
\bibitem{martinez2017} H. Martinez et al., ``Mental workload monitoring: New perspectives from neuroscience,'' \emph{Frontiers in Psychology}, 2017.

% Psychometrics and Validation Theory
\bibitem{messick1989} S. Messick, ``Validity,'' in \emph{Educational Measurement}, R. L. Linn, Ed., 3rd ed. Macmillan, 1989.
\bibitem{cronbach1955} L. J. Cronbach and P. E. Meehl, ``Construct validity in psychological tests,'' \emph{Psychological Bulletin}, 1955.
\bibitem{campbell1959} D. T. Campbell and D. W. Fiske, ``Convergent and discriminant validation by the multitrait-multimethod matrix,'' \emph{Psychological Bulletin}, 1959.

% ML Accountability and Auditing
\bibitem{raji2020} I. D. Raji, M. Mitchell, T. Gebru et al., ``Closing the AI accountability gap: Defining an end-to-end framework for internal algorithmic auditing,'' \emph{Proc. FAT*}, 2020.
\bibitem{mitchell2019} M. Mitchell et al., ``Model cards for model reporting,'' \emph{Proc. FAT*}, 2019.
\bibitem{gebru2021} T. Gebru et al., ``Datasheets for datasets,'' \emph{Comm. ACM}, vol. 64, no. 12, pp. 86--92, 2021.
\bibitem{buolamwini2018} J. Buolamwini and T. Gebru, ``Gender shades: Intersectional accuracy disparities in commercial gender classification,'' \emph{Proc. MLR}, 2018.

% Human Factors and HMS
\bibitem{parasuraman2000} R. Parasuraman et al., ``A model for types and levels of human-machine systems,'' \emph{IEEE Trans. Syst. Man Cybern.}, 2000.
\bibitem{lee2004} J. D. Lee and K. A. See, ``Trust in automation: Designing for appropriate reliance,'' \emph{Human Factors}, vol. 46, no. 1, pp. 50--80, 2004.
\bibitem{wickens2008} C. D. Wickens, ``Multiple resources and performance prediction,'' \emph{Theoretical Issues in Ergon. Sci.}, 2008.
\bibitem{gopher1986} D. Gopher and E. Donchin, ``Workload: An examination of the concept,'' \emph{Handbook of Human Perception and Performance}, 1986.

% Physiological Guidelines and Tools
\bibitem{boucsein2012} W. Boucsein, \emph{Electrodermal Activity}. Springer Science \& Business Media, 2012.
\bibitem{cacioppo2007} J. T. Cacioppo et al., \emph{Handbook of Psychophysiology}. Cambridge University Press, 2007.
\bibitem{neurokit2021} D. Makowski et al., ``NeuroKit2: A Python toolbox for neurophysiological signal processing,'' \emph{Behavior Research Methods}, 2021.

% Datasets
\bibitem{swell2014} S. Koldijk et al., ``The SWELL-KW dataset for research on stress and user modeling,'' \emph{Proc. ICMI}, 2014.
\bibitem{wesad2018} P. Schmidt et al., ``Introducing WESAD, a multimodal dataset for wearable stress and affect detection,'' \emph{Proc. ICMI}, 2018.

% Recent Surveys and Standards
\bibitem{chiossi2024} F. Chiossi et al., ``PhysioCHI: Towards best practices for integrating physiological signals in HCI,'' \emph{Proc. CHI}, 2024.
\bibitem{ieee7000} IEEE Standard 7000-2021, ``Standard model process for addressing ethical concerns during system design,'' IEEE, 2021.

% Methodology and Stats
\bibitem{permutation2010} P. Good, \emph{Permutation, Parametric, and Bootstrap Tests of Hypotheses}. Springer, 2010.
\bibitem{pseudo2015} S. E. Maxwell et al., \emph{Designing Experiments and Analyzing Data}. Routledge, 2015.
\bibitem{zhou2018} Z.-H. Zhou, ``A brief introduction to weakly supervised learning,'' \emph{National Science Review}, vol. 5, no. 1, pp. 44--53, 2018.

\end{thebibliography}

\section*{Author Biographies}
\textbf{Kotaro Furukawa} is currently pursuing the B.S. degree. Their research interests include affective computing, physiological signal processing, and the development of robust human-machine systems. They are the lead developer of the LabelScope audit framework for physiological computing.

\end{document}
